\section{Removing the pole bound}
\label{sec:boundary}

We continue the induction of Theorem~\ref{thm:coefficient-finiteness},
in the lexicographic order on $(n,n-t)$.  At the present pair $(n,t)$,
Proposition~\ref{prop:compact-finiteness} has established the compact
assertion.  At every intermediate pair $(n,s)$ with $t<s<n$, we may use the
full-frame assertion and its perfected-row refinement,
Proposition~\ref{full:perfected-row-characters}.  At $(s,t)$ with
$s<n$, we may use compact finiteness for the admissible coefficients
of Definition~\ref{def:compact:admissible}, including the norm-character
filtrations of Corollary~\ref{cor:compact:norm-filtration}.
The full-frame assertions at $(n,t)$ will be proved in the next
section. On an exact height-$s$ stratum of the original deformation,
$P_n$ is the acting stabilizer and $P_s$ is its commuting auxiliary
connected-frame group. After the relative splitting, the same $P_s$
also acts as the stabilizer of the smaller height-$s$ deformation
disc. Both roles enter the final compact-duality calculation.

We use the algebraic coefficient convention throughout this section.
For example, $B_{\mathrm{full},s}^{\mathrm{perf}}$ denotes the union of
all finite connected and \mbox{\'etale} marking levels for the pair
$(n,s)$ and of their finite Frobenius root extensions.  On a profinite
source, a cochain in this union uses one finite marking and root stage
and one pole bound.  The superscript $\mathrm{perf}$ in this notation
does not include completion.  Powering identifies each fixed root
extension with the original coefficient problem, with all actions and
line bundles transported together.  Thus the uniform full-connected,
fixed-\'etale-level bounds and the fixed scalar subgroup from an earlier
induction stage persist in this root union.

\begin{proposition}[Removal of the pole bound]
\label{prop:pole-removal}
Let $1\leq t<n<p-1$, and retain the notation $A$, $x=x_t$, and $C_1$
of Proposition~\ref{prop:finite-rings}.  The bounded-pole continuous
cohomology
\[
  H^*(P_n,C_1[1/x]\otimes\omega^e),\qquad e\in\Z,
\]
has finite total dimension over $k$, where $\omega$ is the invariant
cotangent line of the universal connected formal group.  Consequently
$H^*(P_n,B_U)$ has finite total dimension at every finite joint frame
level $U$; in particular this holds for every $B_{K,V}$.
\end{proposition}

We prove the proposition after establishing the relative support
calculation.  Powers of $\omega$ are retained in that calculation;
writing the untwisted coefficient in a formula suppresses only that
fixed factor.

\subsection{The height filtration of support}

Put
\[
 J_s=(x_t,\ldots,x_{s-1})\subset A,
 \qquad r=s-t,
 \qquad t<s\leq n,
\]
and use the convention $x_n=1$.  Successive localization triangles for
the closed sets $V(J_s)$ filter the support complex
$R\Gamma_{(x)}(C_1\otimes\omega^e)$ by the terms
\begin{equation}
 \label{eq:boundary-support-rows}
 R\Gamma_{J_s}(C_1\otimes\omega^e)[1/x_s]
 \simeq
 H^{s-t}_{J_s}(C_1\otimes\omega^e)[1/x_s][-(s-t)].
\end{equation}
Here $C_1$ is finite free over $A$, and $x_t,\ldots,x_{s-1}$ is a
regular sequence on this coefficient.  This proves the concentration
in the displayed degree.  The localization triangles can be computed
by finite \v{C}ech complexes, so all denominators in any particular
element or cochain are finite.  Equivalently the local row is the
direct limit of Koszul fractions on nilpotent neighborhoods, with
continuous coefficient splittings obtained by power-series
truncation.  Applying continuous $P_n$-cochains therefore gives the
same finite support filtration.

The last term, $s=n$, needs no relative deformation.  If $L$ is the
finite free coefficient on the full formal disc, the residue pairing
identifies the continuous dual of its top local cohomology with
\[
 \Hom_A(L,\omega_A),
 \qquad
 \omega_A=\bigwedge^{n-t}\Omega^1_{A/k,\mathrm{cts}}.
\]
Continuous duality for the orientable group $P_n$, of dimension $n^2$
(Lemma~\ref{lem:compact:duality}),
then expresses its $P_n$-cohomology as the dual of compact-coefficient
cohomology in complementary degree.  The dual bundle is one of the
allowed compact extras: finite duality for $C_1/A$ and adjunction
identify its canonical line with the volume and conormal lines used
in Proposition~\ref{prop:compact-finiteness}.  That proposition proves
finiteness for $s=n$.  We now treat $t<s<n$.

\subsection{Marked deformations over Artin parameter rings}

On the reduced exact height-$s$ stratum, take a full connected Honda
marking and an integral basis of the surviving \'etale Tate module.
There are $e_s=n-s$ surviving basis vectors.  The marking and basis
towers have group
\[
 P_s\times\GL_{e_s}(\Zp).
\]
Finite \'etale covers lift uniquely across nilpotent thickenings.  We
always lift a finite part of these towers to a fixed nilpotent
neighborhood before making further constructions.

We use the relative marked-deformation theorem over the nilpotent
affine neighborhoods just constructed.  Here is its precise
application.  Let $A_{\mathrm{nil}}\twoheadrightarrow S_0$ be one
such characteristic-$p$ algebra with nilpotent kernel $I$ and reduced
quotient $S_0$, and let
the connected formal $p$-divisible group over $S_0$ be identified
with $\Gamma_s\times_k S_0$ by the full Honda marking.  The
\'etale quotient lifts uniquely across $I$, so its kernel is the
height-$s$ formal $p$-divisible group to which the deformation theorem
is applied.

Work successively over $A_{\mathrm{nil}}/I^a$.  For $a\geq1$ the
next kernel $J=I^a/I^{a+1}$ is square-zero; equip it with the trivial
nilpotent divided powers.  The base has $p=0$, so the display
lifting theorem applies: lifts of the display, with its specified
reduction, are classified by lifts of its Hodge filtration
\cite[Theorem~48 and Corollary~49]{ZinkDisplays}.  In the
dimension-one convention this is the choice of a rank-one Hodge
quotient in a rank-$s$ module.  The difference between two choices
belongs to the Hodge tangent module, a projective module of rank
$s-1$, tensored with $J$.  Since $IJ=0$, the Honda marking
identifies this tangent module on $J$ with $J^{s-1}$.

Lifts of a classifying parameter map from
$k[[v_1,\ldots,v_{s-1}]]$ form a torsor under the same module
$J^{s-1}$.  The Kodaira--Spencer isomorphism for the universal
height-$s$ deformation identifies these two difference actions
\cite[Section~5, pp.~226--227]{Lau2010Tate}.
The universal pullback therefore gives an equivariant bijection
between the torsors of parameter lifts and of marked-deformation
lifts.  Affineness ensures that the required lifts of the finite
projective Hodge quotient have no further gluing obstruction.
The comparison with formal $p$-divisible groups, including the
criterion for lifting morphisms over a nilpotent kernel, is
\cite[Corollary~95]{ZinkDisplays}.  It gives an actual isomorphism
of the pulled-back formal $p$-divisible group with the prescribed
marked deformation. For uniqueness, the nilradical of
$A_{\mathrm{nil}}/I^{a+1}$ is $I/I^{a+1}$ and is nilpotent, because
$S_0$ is reduced. The display functor is therefore fully faithful
by \cite[Proposition~99]{ZinkDisplays}; the injectivity under
nilpotent reduction in \cite[Proposition~40]{ZinkDisplays} gives
uniqueness of the group isomorphism with its prescribed reduction.
These results apply to the connected formal group just specified.
Induction on $a$
proves the classification and its functorial compatibility over
the fixed nilpotent neighborhood.

The successive height congruences set $v_1,\ldots,v_{t-1}$ equal
to zero.  Thus the relative parameter ring used here is
\begin{equation}
 \label{eq:boundary-relative-base}
 R'=k[[v_t,\ldots,v_{s-1}]],
 \qquad \mathfrak m'=(v_t,\ldots,v_{s-1}),
\end{equation}
and $\mathcal G$ denotes the corresponding algebraic
$p$-divisible group, formed from its finite kernels.  The
triangular conormal computation in
Lemma~\ref{lem:artin-splitting} will identify this parameter
filtration with the original transverse filtration.

The finite-stage assertion is coefficientwise.  At a fixed
nilpotence order the parameter values and every specified finite
list of coefficients of the universal isomorphism are elements
of the algebraic marking union, hence descend to one common
finite marking stage.  Equivalently one can record these
coefficients on a sufficiently high finite flat torsion group
and use finite presentation.  This assertion does not select
one finite stage for the whole infinite isomorphism series.
For every particular finite diagram, the Hodge lifts and their
comparison maps involve finite projective modules and finitely
many algebraic operations at that stage.  They are compatible
with the commuting actions by uniqueness.  The explicit
truncation and idempotent sections in
Lemma~\ref{lem:artin-splitting} then supply continuity and a
common bounded pole loss for compact cochain sources.

\begin{lemma}[Relative Artin splitting]
\label{lem:artin-splitting}
Set $R_N=R'/\mathfrak m'^{N}$.  The marked height-$s$ nilpotent
neighborhoods above are flat over $R_N$.  After adjoining compatible
lifts of the surviving \'etale basis, let $E_N$ be their algebraic
union, formed first without the extra coefficient $C_1$.  Then
\[
 E_N/\mathfrak m'E_N=B_{\mathrm{full},s}^{\mathrm{perf}}=:S_s
\]
is perfect, and there is a canonical isomorphism
\begin{equation}
 \label{eq:boundary-artin-product}
 R_N\otimes_k S_s \xrightarrow{\ \sim\ } E_N.
\end{equation}
It is compatible with restriction in $N$, with $P_n$, $P_s$, and the
surviving matrix-group action, and with bounded-pole continuous
cochains on every compact profinite source.  The group translating
the lifts remains the group over $R'$ determined by $\mathcal G$.
\end{lemma}

\begin{proof}
We first verify the parameter assertion, including flatness.  Write
$h$ for the universal isomorphism, oriented by
\[
 h\circ[p]_{\mathrm{original}}=[p]_{\mathcal G}\circ h,
 \qquad c=h'(0)\in E_N^\times.
\]
Successive height-coordinate congruences show that the parameters
below $t$ vanish.  Inductively compare the coefficient of $X^{p^i}$
modulo the earlier height parameters, for $t\leq i<s$.  On that
quotient the two $p$-series start with $x_iX^{p^i}$ and
$v_iX^{p^i}$, respectively.  The isomorphism identity therefore
gives $cx_i=c^{p^i}v_i$ modulo the earlier parameters.  Induction
also identifies the ideals generated by those earlier parameters
on the two sides.  Passing to the conormal module gives
\[
 x_i\equiv c^{p^i-1}v_i+\sum_{j<i}b_{ij}v_j\pmod{J_s^2}.
\]
Changing the orientation of $h$ inverts this matrix.  In particular
the diagonal entries are units.  The $v_i$ and $x_i$ consequently
generate the same transverse ideal and the same ideal-power
filtration on every nilpotent neighborhood.

This conormal calculation is accompanied by a calculation of every
parameter graded piece.  Before marking, regularity of the transverse
parameters gives
\[
 \operatorname{gr}_{J_s}(A[1/x_s]/J_s^N)
  \simeq (A/J_s)[1/x_s][X_t,\ldots,X_{s-1}]/(X)^N.
\]
Finite \'etale marking covers are flat, and this equality therefore
base changes to their special-fiber algebras.  The triangular change
just computed identifies it with the associated graded for $R_N$.
Choose a $k$-linear lift of the special-fiber algebra into the marked
algebra.  Multiplication extends that lift to an $R_N$-linear map
from its tensor product with $R_N$.  The map is an isomorphism on
every associated graded piece, hence is an isomorphism of modules
because the filtration is finite.  This proves flatness over $R_N$;
it is valid regardless of the dimension of the special-fiber algebra
as a $k$-vector space.  Adjoining a finite set of basis lifts is a
finite locally free torsor, so it preserves this flatness.  Its
filtered union does so as well.

The group of compatible lifts is
\begin{equation}
 \label{eq:boundary-relative-translations}
 T'=(T_p\mathcal G)^{e_s}.
\end{equation}
On the closed parameter fiber, Proposition~\ref{prop:frame-torsor}
identifies its finite lift levels with the successive root rings of
the full height-$s$ marked coefficient.  Their union is $S_s$, which
is perfect.

We recall the elementary splitting fact that now applies.  Let $R$
be a finite local Artin $k$-algebra, let $E$ be a flat $R$-algebra,
and suppose $S=E/\mathfrak m_RE$ is perfect.  Choose $q=p^a$ large
enough to kill the nilpotent ideal by $q$th powers and to act
identically on $k$.  For $z\in S$ define
\[
 \sigma(z)=\widetilde{z^{1/q}}^{\,q}\in E,
\]
where the tilde denotes any lift.  The result is independent of the
lift because the difference of two lifts has zero $q$th power.
Using a further root proves independence of sufficiently large $q$.
The characteristic-$p$ power identities show that $\sigma$ is
additive, multiplicative, and $k$-linear.  It is functorial in $E$.
Flatness gives
\[
 \operatorname{gr}_{\mathfrak m_R} E
   =\operatorname{gr}_{\mathfrak m_R}R\otimes_k S.
\]
Thus $r\otimes z\mapsto r\sigma(z)$ is an isomorphism on associated
graded modules and therefore an isomorphism of algebras.  Applying
this fact to $E_N/R_N$ proves
\eqref{eq:boundary-artin-product}.  Functoriality proves all asserted
equivariances and compatibility in $N$.  In particular $P_n$ acts
only on $S_s$ in this description, whereas $P_s$ acts on both factors
in the universal deformation convention.

We check the topological assertion at the same finite stages.
Expansion in the transverse $x_i$ gives continuous additive sections
of their finite truncations.  At a finite \'etale marking stage the
coefficient module is finite projective.  More explicitly, if its
presentation is $e_N(A[1/x_s]/J_s^N)^d$ and $e_0$ is the reduced
idempotent, let $\tau$ lift each coordinate by transverse truncation.
Then $z\mapsto e_N\tau(z)$ is a continuous additive section on
$\operatorname{im}(e_0)$.  Monomial division in the coordinates,
followed by $e_N$, extracts the parameter layers.  A finite
lift-torsor algebra
has the finite free torsion-coordinate bases of
Proposition~\ref{prop:finite-rings}, and the same construction applies
in that basis.  The finitely many structure constants and the
inverse triangular conormal matrix have a common denominator after
localizing at $x_s$.  It follows that these additive lifts and
successive parameter-coefficient extractions have bounded pole loss
and are continuous on a cleared-denominator power-series lattice.
For a compact cochain source all of these choices use one finite
stage and one bound.  In the formula for $\sigma$, extracting the
$q$th root means moving to one further finite root stage; taking
powers and products preserves the stated continuity and merely
changes the finite pole bound.  The inverse of
\eqref{eq:boundary-artin-product} extracts the finitely many parameter
layers by the same maps.  This proves the assertion for cochains.
All tensor products used here have the finite-dimensional algebra
$R_N$ as one factor; no completion or infinite parameter tensor
product has been introduced.
\end{proof}

\subsection{The full-section coefficient across nilpotents}

The remaining coefficient $C_1$ must also be identified on these
thickenings.  Let
\[
 H_c=\mathcal G[p]/\ker(F^t),
 \qquad \operatorname{rank}(H_c)=p^{s-t},
\]
where the quotient is the Frobenius-divided kernel over $R'$.
It is a finite locally free algebraic group scheme.  We write
$C_1^{(s,t)}$ for the smaller labeled ring of
Proposition~\ref{prop:finite-rings}, with total height $s$ and starting
height $t$.

\begin{lemma}[Full sections on a split kernel]
\label{lem:full-sections}
After the lift ascent of Lemma~\ref{lem:artin-splitting}, the pullback
of $C_1$ is a finite disjoint union of coefficients pulled back from
$R'$ of the form
\begin{equation}
 \label{eq:boundary-M}
 M=C_1^{(s,t)}\otimes_{R'}
       \mathcal O(H_c)^{\otimes e_s}.
\end{equation}
An invariant-line twist is the corresponding twist from
$\mathcal G$.  The algebra $M$ is finite free over $R'$, and is free
over $C_1^{(s,t)}$ with finite torsion-coordinate factors.  The
decomposition and its group actions hold on arbitrary nilpotent
test rings.  A sufficiently small surviving matrix congruence
subgroup fixes the components and acts trivially on the factors in
\eqref{eq:boundary-M}.  Translation on these factors uses only a
finite torsion level.
\end{lemma}

\begin{proof}
Recall first the integral full-section presentation.  If $H$ is a
finite monogenic scheme of rank $p^m$, a label homomorphism
$\ell:\Fp^m\to H$ is full when the characteristic polynomial of its
coordinate equals the product over the label sections, with
multiplicities.  Equivalently, for every function $f$ on $H$ after
arbitrary base change,
\begin{equation}
 \label{eq:boundary-norm-sections}
 \det(m_f:\mathcal O(H)\longrightarrow\mathcal O(H))
       =\prod_{v\in\Fp^m} f(\ell(v)).
\end{equation}
The equivalence follows by polynomial substitution in the
characteristic-polynomial identity, or by the resultant formula for
the norm.  Thus this condition is intrinsic and is preserved by
every automorphism of the finite scheme.

For the Frobenius-divided $[p]$-kernel,
Lemma~\ref{rings:full-sections-presentation} identifies this
full-section algebra with $C_1$, including after arbitrary
nilpotent base change.

Over the split lift tower the divided kernel is
\[
 H=H_c\times E_{\mathrm{et}},
 \qquad E_{\mathrm{et}}\simeq\Fp^{e_s}.
\]
The projection of a full tuple gives a surjection
\[
 b:\Fp^{n-t}\twoheadrightarrow\Fp^{e_s}.
\]
Indeed, if a component of $E_{\mathrm{et}}$ were omitted, take the
function which is zero on that component and one on every other
component.  Its norm is zero, but its product over the proposed
sections is one, contradicting \eqref{eq:boundary-norm-sections}.
Projection to the constant \'etale group is locally constant on the
test scheme.  We can therefore decompose the full-section scheme
into the finitely many open and closed components indexed by these
surjections.

Fix $b$, put $K=\ker b$, and choose a complement $W$ to $K$ in the
label space.  Its dimensions are $s-t$ and $e_s$, respectively.  A
label homomorphism with this projection is exactly the data of a
homomorphism $K\to H_c$ and $e_s$ arbitrary points of $H_c$, the
images of a basis of $W$.  On each \'etale component, the proposed
tuple is the translate of the kernel tuple by the associated point
of $H_c$. To test fullness on one component, take an arbitrary
function on that copy of $H_c$ and extend it by $1$ on every other
component. Equation~\eqref{eq:boundary-norm-sections} then reduces
exactly to the norm identity for that one component. Conversely,
the product of the componentwise identities gives the identity for
every function on the whole disjoint union. Translation preserves
these norm identities. Thus the original tuple is full if and only
if its kernel tuple is full for $H_c$, over every test algebra.
The two constructions are inverse without reducing that algebra.
The component is therefore
\[
 \operatorname{Full}(H_c)\times H_c^{e_s},
\]
whose ring is \eqref{eq:boundary-M}.

The argument uses algebraic translations on the finite group
scheme.  It does not evaluate a formal power series on an arbitrary
localized coordinate; the norm identity supplies the required
substitution invariance within the finite algebra.  Finite freeness
follows from the smaller $C_1$ calculation and finite local
freeness of $H_c$.  All operations are functorial for changes of
deformation framing.  Changes of surviving Tate basis act on these
particular data through a finite level, since the tuple and the
divided kernel are finite-level objects.  A sufficiently deep
congruence subgroup fixes the chosen labels, complement, and
connected factors.  The same finite-level observation applies to
translations and to each fixed invariant-line or torsion twist.
\end{proof}

Let $r=s-t$.  For a coefficient $M$ in
\eqref{eq:boundary-M}, with any of the fixed twists just discussed,
put
\[
 \mathcal H(M)=H^r_{\mathfrak m'}(M).
\]
This is a discrete union of finite-dimensional $k$-vector spaces.
For example, it is computed by the quotients
\begin{equation}
 \label{eq:boundary-local-union}
 \mathcal H(M)
 =\colim_{N}
  M/(v_t^N,\ldots,v_{s-1}^N)M,
\end{equation}
whose transition is multiplication by $v_t\cdots v_{s-1}$.
These maps are injective, since $M$ is free over $R'$ and
\[
 ((v_t^{N+1},\ldots,v_{s-1}^{N+1}):v_t\cdots v_{s-1})
    =(v_t^N,\ldots,v_{s-1}^N).
\]
The rectangular ideals here and the powers of $\mathfrak m'$ are
cofinal.  The compatible Artin splittings therefore identify the
full split local row with the ordinary tensor product
\begin{equation}
 \label{eq:boundary-split-local-row}
 S_s\otimes_k\mathcal H(M).
\end{equation}
The action of $P_n$ is on $S_s$ alone.  A fixed finite projective
resolution for $P_n$ and the finite-dimensional stages in
\eqref{eq:boundary-local-union} give, with the specified cochain
convention,
\begin{equation}
 \label{eq:boundary-split-cohomology}
 H^i(P_n,S_s\otimes_k\mathcal H(M))
   =H^i(P_n,S_s)\otimes_k\mathcal H(M).
\end{equation}
The same formulas hold after adjoining any fixed finite list of
torsion-coordinate factors from $\mathcal G$.  In particular the
identification keeps the complete nilpotent local-cohomology term.

\subsection{Forgetting the varying translation torsor}

The earlier full-frame assertion at $(n,s)$ supplies one nontrivial
scalar $a\in1+p\Zp$ acting identically on $H^*(P_n,S_s)$.
After taking a power, it also fixes the finite label and torsion
data in $M$.  It consequently acts identically on
\eqref{eq:boundary-split-cohomology}.  It conjugates $T'$ by
$[a]$ or $[a^{-1}]$, according to the dual-standard convention.
The following calculation treats either convention.

\begin{lemma}[Scalar conjugation and relative translation cohomology]
\label{boundary:translation-descent}
On the cohomology rows in \eqref{eq:boundary-split-cohomology}, some
subgroup $T'_h=[p]^hT'$ acts trivially.  Over an Artin parameter
neighborhood, write
\[
 \Lambda_N=R_N[[D_1,\ldots,D_\delta]],
 \qquad \delta=e_s(s-1),\qquad I=(D_1,\ldots,D_\delta)
\]
for the distribution algebra of $T'_h$, transported along
$T'\simeq T'_h$.  If $V$ is a row with trivial $T'_h$ action, then
\begin{equation}
 \label{eq:boundary-relative-koszul}
 H^j(T'_h,V)
   \simeq
 V\otimes_{R_N}\bigwedge^j(I/I^2)^\vee,
 \qquad 0\leq j\leq\delta.
\end{equation}
This description is functorial under changes of distribution
coordinates and under all the commuting actions.
\end{lemma}

\begin{proof}
Since $s\geq2$, the special Cartier dual of $\mathcal G$ is
connected, of dimension $s-1$.  Its formal completion is smooth over
$R'$.  The completed distribution algebra of $T'$ is therefore a
power-series algebra over $R'$ on $\delta$ parameters.  The
representations under consideration are locally killed by a
parameter power and locally factor through a finite torsion level.
These are the continuous torsion modules for this distribution
algebra.  To check the topology, modulo $\mathfrak m'$ the finite
dual-kernel quotients are cofinal with powers of their augmentation
ideal, because the special dual is connected.  At a fixed Artin
thickness this cofinality lifts by multiplying exponents by a
nilpotence bound.  Conversely the coordinates of $[p]$ have no
linear term in characteristic $p$, so their iterates tend to zero
in the augmentation topology.  The inverse limit of these finite
distribution quotients is precisely $\Lambda_N$.

Suppose first that $a$ conjugates translations by $[a]$, and let
$\sigma_a=[a]^*$ on the distribution algebra of $T'$.  As $a$ acts
identically on the row, the ideal
\[
 J_a=(\sigma_a(f)-f:f\in\Lambda_N)
\]
annihilates it.  Write $F$ for the formal group law on the Cartier
dual.  The identity
\[
 F([a](D),[-1](D))=[a-1](D)
\]
shows that the coordinates of $[a-1](D)$ belong to $J_a$.
Conversely, $F(D,[a-1](D))=[a](D)$ shows that
$[a](D)-D$ belongs to the ideal of those coordinates.  Thus
\[
 J_a=([a-1]_1(D),\ldots,[a-1]_\delta(D)).
\]
If $a-1=p^hu$ with $u\in\Zp^\times$, the automorphism $[u]$ gives
\[
 J_a=[p^h]^*(I).
\]
This is the image of the augmentation ideal for restriction to
$[p]^hT'$, so that subgroup acts trivially.  With $a^{-1}$ in place
of $a$, its difference from $1$ has the same $p$-adic valuation and
the argument is identical.  The chosen scalar is uniform on the
special cohomology and the coefficient $M$ is one fixed finite
construction, so $h$ is uniform over its local-cohomology union.

It remains to justify the relative cohomology calculation.  The
sequence $D_1,\ldots,D_\delta$ is regular in $R_N[[D]]$, even though
$R_N$ itself can have nilpotents: successive quotients are the
power-series rings in the remaining variables over $R_N$, and
multiplication by the next variable is injective.  Its Koszul
complex is a finite free resolution of $R_N$ over $\Lambda_N$.
Applying continuous $\Hom$ to a row annihilated by $I$ gives zero
differentials, and hence \eqref{eq:boundary-relative-koszul}.  No
flatness of the row over $R_N$ is required.

This Ext calculation is also the torsor-descent calculation being
used here.  At finite levels the lift algebras are finite locally
free, with the unit a local basis vector, and their augmented
descent diagrams are split as base-module diagrams on affine
covers.  Their flat union supplies the relative continuous bar
construction.  Over $R_N$, its completed distribution resolution
and the Koszul resolution are resolutions split over the base:
unit insertion and monomial division give continuous contractions
of their augmented complexes.  An induced completed bar term
$\Lambda_N\mathbin{\widehat\otimes}_{R_N}W$ lifts along a
continuously $R_N$-split surjection by its $R_N$-linear splitting;
finite free Koszul terms have the same lifting property.  Induction
along the two resolutions therefore gives continuous comparison
maps in both directions and comparison homotopies.
Alternatively continuous $\Hom$ from each finite free Koszul term
commutes with the union of torsion coefficients.  Consequently
both resolutions compute the same relative Ext.  Its exterior
description is canonically the exterior algebra on the dual
augmentation conormal, so nonlinear changes of parameters give
the asserted action.  All the calculations are made at a common
Artin thickness before passing to the local-cohomology union.
\end{proof}

The exterior bundles in \eqref{eq:boundary-relative-koszul} are
tensor constructions from the Lie bundle of $\mathcal G^\vee$, with
the standard matrix-label action.  A small matrix congruence
subgroup acts trivially on that label representation in
characteristic $p$.  They are therefore among the extra bundles
allowed in the smaller compact calculation.  Notice that this
argument uses the dimension $s-1$ of the relative Cartier dual;
it does not replace the varying translation group by a constant
formal group.

\subsection{Individual matrix-group rows}

We need finiteness after taking the surviving matrix invariants on
individual $P_n$-cohomology rows.  Here is the algebraic implication
which supplies it from the full-connected, fixed-\'etale-level
assertion.

\begin{lemma}
\label{boundary:row-finiteness}
Let $V$ be a torsion-free pro-$p$ compact $p$-adic analytic group,
and let $C$ be a bounded complex of smooth $k[V]$-modules.  If
$R\Gamma(V,C)$ has finite-dimensional cohomology in each degree,
then
\[
 H^a(V,H^i(C))
\]
is finite-dimensional for every $a,i$.  These groups are zero
outside finitely many $(a,i)$ when $C$ has bounded cohomology.
\end{lemma}

\begin{proof}
Take the full $k$-linear dual, with its product topology, and put
$\Omega=k[[V]]$.  This duality is exact from discrete vector spaces
to pseudocompact vector spaces and exchanges smooth $V$-modules
with pseudocompact $\Omega$-modules.  The ring $\Omega$ is local and
Noetherian.  The dual bounded complex $D=C^\vee$ has
degreewise finite derived completed reduction
\[
 k\mathbin{\widehat\otimes}^{\mathbb L}_{\Omega}D
       \simeq R\Gamma(V,C)^\vee.
\]
We recall why this implies finite generation of each $H^j(D)$.
All differentials in $D$ have closed image, so its cohomology is
pseudocompact. At the highest nonzero cohomological degree $b$,
right t-exactness of derived completed tensor gives
\[
 H^b\!\left(k\mathbin{\widehat\otimes}^{\mathbb L}_\Omega D\right)
     =k\mathbin{\widehat\otimes}_\Omega H^b(D).
\]
Thus the ordinary completed reduction of $H^b(D)$ is finite.
Lift a finite set of
generators of that reduction to cycles.  Their $\Omega$-span is
closed, because it is the image of a finite free pseudocompact
module.  Its quotient has zero completed reduction, and is zero
by compact Nakayama: a nonzero pseudocompact quotient has a
nonzero finite simple quotient, which for the local algebra
$\Omega$ is $k$.  Thus these finitely many cycles generate
$H^b(D)$.

Map the finite free module on these cycles into $D$, in degree
$b$, and take its cone.  The cone has no cohomology in degree
$b$ and still has degreewise finite derived reduction.  Repeat
downwards.  In recovering the cohomology of the preceding
complexes, all kernels inside finite free modules are finitely
generated by Noetherianity.  The exact sequences of these cones
therefore prove finite generation of every original $H^j(D)$;
only finitely many cone steps are needed for each degree.
For a finitely generated pseudocompact $\Omega$-module $M$, choose
a resolution of $M$ whose terms are finite-rank free
pseudocompact $\Omega$-modules.  Such a resolution is constructed
successively: the kernel at each step is closed and finitely
generated because $\Omega$ is Noetherian.  Completed tensoring
this resolution with $k$ gives finite-dimensional terms in every
degree.  Consequently every
$\operatorname{Tor}^{\Omega}_a(k,M)$ is finite-dimensional.
The finite projective dimension of the trivial module $k$
supplies the vanishing range $a>\dim(V)$.  Applying these facts
to the finitely generated cohomology modules of $D$ and using
duality proves the assertion for $H^a(V,H^i(C))$.  The bounded
range for $i$ gives the asserted finite total range.
\end{proof}

Apply the lemma to
$C=R\Gamma(P_n,B_{\mathrm{full},s}^{\mathrm{perf}})$, using the
fixed finite $P_n$-resolution.  Its terms and cohomology are smooth
for the surviving matrix group: a cochain uses one finite stage.
The finite group resolution computes the same invariants in the
smooth category and with the bounded-pole convention, since in
either case its terms are the same finite sums and retracts of
the coefficient modules.
Descent of the full \'etale frame tower gives
\[
 R\Gamma(V,C)
    \simeq R\Gamma(P_n,B_{\mathrm{conn},V,s}^{\mathrm{perf}}).
\]
At finite levels this is ordinary finite \'etale torsor descent;
the augmented coefficient diagrams have module splittings, and
the common-stage convention passes the assertion to the union.
The earlier full-frame assertion makes the right-hand side finite.
Thus, for all sufficiently small open matrix groups $V$,
\begin{equation}
 \label{eq:boundary-row-finiteness}
 \dim_k H^a\!\left(V,
       H^i(P_n,B_{\mathrm{full},s}^{\mathrm{perf}})\right)<\infty.
\end{equation}
An arbitrary compact open follows by finite-index descent.  The
commuting $P_s$-action on these groups is smooth and therefore
factors through a finite quotient.  We also need the
representation-theoretic refinement furnished by the completed
earlier stage.  Put
\[
 \nu_s:P_s\xrightarrow{\operatorname{Nrd}}\Zp^\times
             \longrightarrow\Fp^\times\subset k^\times .
\]
By Proposition~\ref{full:perfected-row-characters}, each actual
row
\begin{equation}
 \label{eq:boundary-norm-row}
 Q_{a,i,V}
   =H^a\!\left(V,H^i(P_n,B_{\mathrm{full},s}^{\mathrm{perf}})\right)
\end{equation}
has a finite $P_s$-stable filtration with one-dimensional
successive quotients $k(\nu_s^j)$.  The exponents may depend on
the row.  Its linear dual has the dual filtration, with
successive quotients $k(\nu_s^{-j})$.  This is the additional
structure that permits these particular constant factors in
the smaller-disc compact calculation.

Only the previously proved induction stage $(n,s)$, with $s>t$,
is used here.
After the full-frame calculation and
Proposition~\ref{full:perfected-row-characters} are proved at
$(n,t)$, the same conclusions become available for its use in
subsequent pole-removal stages.  Thus both the finite-dimensional
row assertion and its norm-character refinement follow the
same lexicographic induction.

\subsection{Final descent and compact duality on the smaller disc}

\begin{proof}[Proof of Proposition~\ref{prop:pole-removal}]
For an intermediate support row $t<s<n$, perform the marked and
split ascent of Lemmas~\ref{lem:artin-splitting} and
\ref{lem:full-sections}.  The resulting coefficient is
\eqref{eq:boundary-split-local-row}.  First take $P_n$-cohomology,
then forget the subgroup $T'_h$ supplied by
Lemma~\ref{boundary:translation-descent}.  The cohomology rows are
the tensors in \eqref{eq:boundary-split-cohomology} with the finite
exterior bundles in \eqref{eq:boundary-relative-koszul}.

Take next the invariants of a sufficiently small surviving matrix
group $V$, chosen also to fix all finite label and torsion factors
being considered.  The finite resolution for $V$ commutes with the
discrete union \eqref{eq:boundary-local-union}.  Its individual
cohomology terms consequently have the form
\begin{equation}
 \label{eq:boundary-before-Ps}
 Q\otimes_k H^r_{\mathfrak m'}(M\otimes_{R'}F).
\end{equation}
Here $Q$ is one of the actual rows
\eqref{eq:boundary-norm-row}, and $F$ is a finite free tensor
construction from the Lie bundle of $\mathcal G^\vee$,
invariant lines, and the fixed additional finite
torsion-coordinate bundles.  In particular $Q$ is finite and
its $P_s$-composition factors are norm powers.  Every fixed
bar degree adds only such universal torsion-coordinate factors
to $F$; the special coefficient producing $Q$ remains
$B_{\mathrm{full},s}^{\mathrm{perf}}$.

We finally take $P_s$-cohomology.  The residue pairing on the
$r$-dimensional disc identifies the continuous dual of the local
term in \eqref{eq:boundary-before-Ps} with
\[
 L_{M,F}=\Hom_{R'}(M\otimes_{R'}F,\omega_{R'}),
 \qquad
 \omega_{R'}=\bigwedge^r\Omega^1_{R'/k,\mathrm{cts}}.
\]
The group $P_s$ is orientable of dimension $s^2$, and continuous
group duality therefore gives
\begin{equation}
 \label{eq:boundary-Matlis-duality}
 H^b\!\left(P_s,
    Q\otimes_kH^r_{\mathfrak m'}(M\otimes F)\right)^\vee
 \simeq
 H^{s^2-b}(P_s,Q^\vee\otimes_kL_{M,F}).
\end{equation}
First we verify that $L_{M,F}$ is an admissible compact
coefficient at the smaller pair $(s,t)$, whose total height is
strictly less than $n$.  The algebra $M$ is free over
$C_1^{(s,t)}$ and consists of finite torsion-coordinate factors.
On each successive labeled boundary stratum of that smaller
disc, finite root transport identifies the construction with
the corresponding height-stratum construction.  After finite
connected and \'etale markings and a finite splitting level,
its divided connected kernels and finite \'etale sections have
the required constant frames.  The Lie and invariant-line
factors have the same property.  Finite tensor constructions
and linear duals preserve these filtrations.  If a
multiplicative connected kernel occurs, its character factors
are summands of the finite regular translation representation;
its associated coefficient is the finite root ring allowed in
the compact assertion.  Finally finite duality gives
\[
 \Hom_{R'}(M\otimes F,\omega_{R'})
 \simeq
 \Hom_{C_1^{(s,t)}}(M\otimes F,\omega_{C_1^{(s,t)}}),
\]
and the volume construction and conormal adjunction supply the
required filtration for this canonical-line twist.  This proves
admissibility of $L_{M,F}$.

Now use the finite $P_s$-stable filtration of $Q^\vee$ from
\eqref{eq:boundary-norm-row}.  Tensoring it with $L_{M,F}$
gives a finite filtration with successive coefficients
\[
 L_{M,F}\otimes_k k(\nu_s^{-j}).
\]
Lemma~\ref{lem:compact:norm-line} identifies the constant norm
line through the determinant of the universal height-$s$
$p$-divisible group and proves its admissibility on every
required stratum.  Thus these successive coefficients are
admissible by Corollary~\ref{cor:compact:norm-filtration}.
Their compact $P_s$-cohomology is finite by the earlier compact
assertion at $(s,t)$.  The finite-dimensional filtration of
$Q^\vee$ splits over $k$ as a filtration of vector spaces;
after tensoring with $L_{M,F}$, its coefficient sequences
therefore have continuous underlying module splittings.
Their long exact cohomology sequences prove finiteness for
$Q^\vee\otimes_kL_{M,F}$.  Consequently
\eqref{eq:boundary-Matlis-duality} is finite-dimensional.
The acting group in this application is $P_s$ throughout.

These successive computations are legitimate descent spectral
sequences for the coefficient under consideration.  At an Artin
stage, adjoining finite free translation coordinates commutes with
$P_n$-cochains because those coordinates come from $\mathcal G/R'$
and $P_n$ fixes $R'$.  Flat torsor descent and the finite-module
splittings described above give the bar construction before taking
the local row.  The compatible finite-thickness calculations then
give it on that row.  Thus one may filter first by $P_n$-, then by
translation-, and then by matrix-group cohomology before taking
$P_s$-cohomology.  Each final contribution is a group of the form
\eqref{eq:boundary-Matlis-duality}.  All group degrees are
nonnegative, and the one local-cohomology shift is the fixed
integer $r$.

There is still the finite translation quotient $T'/T'_h$ to
forget.  Use its augmented bar construction.  In bar degree $q$
it adds $q$ finite torsion-coordinate factors of $\mathcal G$,
of one fixed finite torsion order.  They are finite free over
$R'$, and $T'_h$ acts trivially on them.  The calculation just
given applies in every bar degree, with these factors included
in $M\otimes F$.  A further matrix congruence restriction fixes
their finite-level label actions.
For every resulting open matrix group, the representation
assertion of Proposition~\ref{full:perfected-row-characters}
still applies to its special cohomology rows.  The factors
added in the finite quotient bars remain in the universal
bundle $M\otimes F$, so the same norm-filtration argument
applies in each bar degree.
There are only finitely many
bar degrees in a fixed total degree, so this descent preserves
finite-dimensionality in each degree.  The same reasoning forgets
the finite quotients used to fix label components or to replace a
matrix group by a small open subgroup.  Forgetting the full
connected marking is exactly the $P_s$ step above.  We have proved
that every support row in \eqref{eq:boundary-support-rows} has
finite-dimensional $P_n$-cohomology.  Its intrinsic $P_n$
cohomological bound is $n^2$, so its total cohomology is finite.

The endpoint $s=n$ was treated by compact duality at the beginning
of the section.  The finite height-support filtration therefore
has finite total $P_n$-cohomology.  Compact finiteness on $C_1$,
together with the localization triangle
\[
 R\Gamma_{(x)}(C_1\otimes\omega^e)
 \longrightarrow C_1\otimes\omega^e
 \longrightarrow C_1[1/x]\otimes\omega^e,
\]
proves the first assertion of the proposition.

For completeness, this yields every finite frame level as follows.
After the first \'etale basis level, the leading connected marking
is tame and decomposes into powers of the invariant line.  All
subsequent sufficiently deep normal finite joint levels have
finite $p$-group deck transformations.  Their regular
representations over $k$ have finite filtrations by trivial
representations.  The associated coefficient bundles consequently
have finite filtrations by the invariant-line coefficients already
treated.  The filtrations split as modules on localization, so
they give exact sequences of bounded-pole continuous cochains.
This proves finite total cohomology at all such sufficiently deep
levels.  Any prescribed finite joint level is below one of these.
Finite torsor descent and its cohomological spectral sequence
prove finite-dimensionality in each degree there as well.  The
fixed cohomological bound for $P_n$ again gives finite total
dimension.  Every step admits arbitrary compact profinite-source
parameters by the common-thickness, common-stage, and bounded-pole
constructions above.  The asserted finiteness is finiteness of the
resulting cohomology groups over $k$.
\end{proof}
